% ═══════════════════════════════════════════════════════════════════════════════
%  Soutenance — Prédiction de fin d'orage (Météorage / Data Battle IA PAU 2026)
% ═══════════════════════════════════════════════════════════════════════════════
\documentclass[aspectratio=169,10pt]{beamer}

% ── Thème & Couleurs ─────────────────────────────────────────────────────────
\usetheme{metropolis}
\usepackage[french]{babel}
\usepackage[utf8]{inputenc}
\usepackage[T1]{fontenc}

\usepackage{graphicx}
\usepackage{booktabs}
\usepackage{tabularx}
\usepackage{amsmath,amssymb}
\usepackage{xcolor}
\usepackage{colortbl}
\usepackage{tikz}
\usepackage{appendixnumberbeamer}
\usepackage{caption}
\captionsetup{font=tiny,labelfont=bf}

% ── Palette Météorage ────────────────────────────────────────────────────────
\definecolor{meteoBlue}{HTML}{1D3557}
\definecolor{meteoRed}{HTML}{E63946}
\definecolor{meteoGray}{HTML}{F1FAEE}
\definecolor{meteoCyan}{HTML}{457B9D}
\definecolor{meteoGold}{HTML}{E9C46A}
\definecolor{meteoGreen}{HTML}{2A9D8F}

\setbeamercolor{palette primary}{bg=meteoBlue,fg=white}
\setbeamercolor{frametitle}{bg=meteoBlue,fg=white}
\setbeamercolor{title separator}{fg=meteoRed}
\setbeamercolor{progress bar}{fg=meteoRed,bg=meteoBlue!20}
\setbeamercolor{alerted text}{fg=meteoRed}
\setbeamercolor{block title}{bg=meteoCyan,fg=white}
\setbeamercolor{block body}{bg=meteoGray}

\setbeamertemplate{navigation symbols}{}
\metroset{sectionpage=progressbar}

% ── Macros ────────────────────────────────────────────────────────────────────
\newcommand{\hig}[1]{\textcolor{meteoRed}{\textbf{#1}}}
\newcommand{\cmark}{\textcolor{meteoGreen}{\checkmark}}
\newcommand{\xmark}{\textcolor{meteoRed}{\texttimes}}

% ═════════════════════════════════════════════════════════════════════════════
\title{Prédiction de fin d'orage}
\subtitle{Data Battle IA PAU 2026 — Météorage}
\date{Mars 2026}
\author{Colin Geindre \and Alexandre Hayoun \and Tao Levray}
\institute{Équipe : \textbf{Les Vieux Briscards}}

% ═════════════════════════════════════════════════════════════════════════════
\begin{document}

% ── Page de titre ─────────────────────────────────────────────────────────────
\maketitle

% ── Plan ──────────────────────────────────────────────────────────────────────
\begin{frame}{Plan}
\tableofcontents
\end{frame}


% ╔═══════════════════════════════════════════════════════════════════════════╗
% ║  PARTIE 1 — CONTEXTE & EDA                                              ║
% ╚═══════════════════════════════════════════════════════════════════════════╝
\section{Phase 1 — Exploration des données}

% ── Contexte ──────────────────────────────────────────────────────────────────
\begin{frame}{Contexte et problématique}
\begin{columns}[T]
\begin{column}{0.55\textwidth}
\textbf{Problème métier :}
\begin{itemize}
  \item Les alertes orage sont maintenues \hig{30\,min} après le dernier éclair CG
  \item Approche conservatrice $\Rightarrow$ pertes d'exploitation aéroportuaires
\end{itemize}

\vspace{0.8em}
\textbf{Objectif :}
\begin{itemize}
  \item Estimer en \textbf{temps réel} le temps résiduel avant la fin de l'orage
  \item Réduire les faux maintiens d'alerte tout en garantissant la sécurité
\end{itemize}
\end{column}
\begin{column}{0.42\textwidth}
\begin{block}{Dataset}
\small
\begin{tabular}{@{}rl@{}}
Éclairs  & 507\,071 \\
Aéroports & 5 \\
Sessions  & 769 \\
Éclairs / alerte & 56\,599 \\
\end{tabular}
\end{block}

\vspace{0.5em}
\small\textit{5 sites : Ajaccio, Bastia, Biarritz, Nantes, Pise}
\end{column}
\end{columns}
\end{frame}


% ── Répartition aéroports ────────────────────────────────────────────────────
\begin{frame}{Répartition par aéroport}
\begin{columns}[c]
\begin{column}{0.6\textwidth}
\includegraphics[width=\linewidth]{figures/fig1_airport_counts.pdf}
\end{column}
\begin{column}{0.38\textwidth}
\begin{itemize}
  \item \textbf{Pise} domine ($>30\%$)
  \item Bastia $\approx 2\times$ Ajaccio malgré la proximité géographique
  \item Nantes : activité la plus faible
\end{itemize}
\end{column}
\end{columns}
\end{frame}


% ── Temporel ──────────────────────────────────────────────────────────────────
\begin{frame}{Patterns temporels}
\begin{columns}[c]
\begin{column}{0.52\textwidth}
\includegraphics[width=\linewidth]{figures/fig2_temporal.pdf}
\end{column}
\begin{column}{0.46\textwidth}
\includegraphics[width=\linewidth]{figures/fig3_heatmap_month_hour.pdf}
\end{column}
\end{columns}
\vspace{0.3em}
\centering\small
\textbf{Saisonnalité estivale} marquée (pic juin--sept.) · \textbf{Cycle diurne} avec maximum 12h--15h UTC
\end{frame}


% ── Aéroport × mois ──────────────────────────────────────────────────────────
\begin{frame}{Profil saisonnier par aéroport}
\centering
\includegraphics[width=0.78\linewidth]{figures/fig4_heatmap_airport_month.pdf}

\vspace{0.3em}
\small
Biarritz : pic \textbf{précoce} (mai--août) · Pise : pic \textbf{tardif} (août--sept.) · Corse : pic \textbf{automnal} (août--oct.)
\end{frame}




% ── Rose des azimuts ──────────────────────────────────────────────────────────
\begin{frame}{Direction des orages (rose des azimuts)}
\centering
\includegraphics[width=0.7\linewidth]{figures/fig6_azimuth_rose.pdf}

\vspace{0.3em}
\small
Chaque aéroport a un \textbf{profil directionnel unique} $\Rightarrow$ pertinence d'une composante spatiale
\end{frame}


% ── Géographie ────────────────────────────────────────────────────────────────
\begin{frame}{Répartition géographique des éclairs}
\centering
\includegraphics[width=0.85\linewidth]{figures/fig7_geo_scatter.pdf}

\vspace{0.3em}
\small
Clusters liés au relief (Corse, Apennins) — distributions homogènes en plaine (Nantes)
\end{frame}




% ── Synthèse EDA ─────────────────────────────────────────────────────────────
\begin{frame}{Synthèse de l'EDA}
\begin{block}{Enseignements clés}
\begin{itemize}
  \item Forte \textbf{saisonnalité} et \textbf{cycle diurne} $\Rightarrow$ features temporelles cycliques
  \item \textbf{Hétérogénéité} entre aéroports $\Rightarrow$ encodage du site
  \item \textbf{Déséquilibre de classes} par construction $\Rightarrow$ formulation régression/survie
  \item \textbf{Structure séquentielle} = information la plus riche $\Rightarrow$ processus ponctuels, RNN
  \item Signaux de \textbf{dissipation} : ratio IC $\nearrow$, intervalles inter-CG $\nearrow$
\end{itemize}
\end{block}
\end{frame}


% ╔═══════════════════════════════════════════════════════════════════════════╗
% ║  PARTIE 2 — COMPARAISON DE MODÈLES                                      ║
% ╚═══════════════════════════════════════════════════════════════════════════╝
\section{Phase 2 — Comparaison de modèles}


% ── Vue d'ensemble ────────────────────────────────────────────────────────────
\begin{frame}{Résultats globaux}
\centering
\includegraphics[width=0.85\linewidth]{figures/comparison_global.pdf}

\vspace{0.3em}
\begin{tabular}{lccccc}
\toprule
\textbf{Modèle} & \textbf{MAE} & \textbf{RMSE} & \textbf{Méd.} & \textbf{Biais} & \textbf{P90} \\
\midrule
Baseline 30 min       & 60.26 & 94.05  & 33.32 & $-$52.32 & 132.48 \\
XGBoost AFT           & 57.28 & 90.69  & 33.18 & $-$26.74 & 149.82 \\
Hawkes Classique      & 83.99 & 124.37 & 56.88 & $-$83.82 & 178.96 \\
\rowcolor{meteoGray}\textbf{Neural Hawkes (GRU)} & \textbf{22.11} & \textbf{32.50} & \textbf{18.00} & \textbf{$-$1.61} & \textbf{40.12} \\
BNN MC Dropout        & 68.20 & 102.98 & 34.31 & +1.42 & 164.40 \\
\bottomrule
\end{tabular}
\end{frame}


% ── XGBoost ──────────────────────────────────────────────────────────────────
\begin{frame}{XGBoost AFT — modèle tabulaire}
\begin{columns}[c]
\begin{column}{0.55\textwidth}
\includegraphics[width=\linewidth]{figures/xgboost_importance.pdf}
\end{column}
\begin{column}{0.43\textwidth}
\begin{itemize}
  \item Modélise $\ln(T)$ via Accelerated Failure Time
  \item 20 features (temporelles, spatiales, physiques)
  \item MAE = 57.28 min ($\sim$ baseline)
  \item \textbf{Limite} : traite chaque éclair \emph{indépendamment}
  \item \texttt{n\_lightnings\_15m} domine
\end{itemize}
\end{column}
\end{columns}
\end{frame}


% ── Hawkes Classique ──────────────────────────────────────────────────────────
\begin{frame}{Processus de Hawkes classique}
\begin{columns}[c]
\begin{column}{0.55\textwidth}
\includegraphics[width=\linewidth]{figures/hawkes_intensity.pdf}
\end{column}
\begin{column}{0.43\textwidth}
\textbf{Intensité conditionnelle :}
$$\lambda(t) = \mu + \sum_{t_i < t} \alpha\beta\,e^{-\beta(t-t_i)}$$

\vspace{0.5em}
\begin{itemize}
  \item MAE = 83.99 min (\hig{pire} modèle)
  \item Noyau exponentiel trop simpliste
  \item \textbf{Mais} : le cadre théorique est validé par le Neural Hawkes
\end{itemize}
\end{column}
\end{columns}
\end{frame}


% ── Neural Hawkes ─────────────────────────────────────────────────────────────
\begin{frame}{Neural Hawkes Process (GRU) — \textcolor{meteoGreen}{meilleur modèle}}
\begin{columns}[T]
\begin{column}{0.55\textwidth}
\textbf{Principe :}\\
Remplacer le noyau exponentiel par un \textbf{GRU} qui apprend la dynamique :
$$\lambda(t) = \text{softplus}\big(W_h \cdot h_t + w_\Delta \cdot \Delta t + b\big)$$

\vspace{0.5em}
\textbf{Avantages :}
\begin{itemize}
  \item Interactions \textbf{non-linéaires} entre features
  \item Dépendances \textbf{à long terme}
  \item Capture la physique de chaque éclair
\end{itemize}
\end{column}
\begin{column}{0.42\textwidth}
\begin{block}{Métriques}
\begin{tabular}{@{}rl@{}}
MAE     & \hig{22.11 min} \\
RMSE    & 32.50 min \\
Médiane & 18.00 min \\
Biais   & $-$1.61 min \\
P90     & 40.12 min \\
\end{tabular}
\end{block}

\vspace{0.5em}
\small $3\times$ meilleur que la baseline\\
Biais \textbf{quasi nul}
\end{column}
\end{columns}
\end{frame}


% ── Distribution des erreurs ──────────────────────────────────────────────────
\begin{frame}{Distribution des erreurs}
\centering
\includegraphics[width=0.85\linewidth]{figures/error_distributions.pdf}

\vspace{0.3em}
\small
Le Neural Hawkes (GRU) a une distribution \textbf{nettement plus resserrée} vers zéro
\end{frame}




% ── Synthèse Phase 2 ────────────────────────────────────────────────────────
\begin{frame}{Synthèse Phase 2}
\begin{block}{Classement final}
\begin{enumerate}
  \item \textbf{Neural Hawkes (GRU)} — MAE = 22.11, biais $\approx$ 0 \cmark
  \item \textbf{XGBoost AFT} — MAE = 57.28, bon modèle tabulaire
  \item \textbf{Baseline} — MAE = 60.26, référence métier
  \item \textbf{BNN MC Dropout} — Incertitude utile mais MAE élevée
  \item \textbf{Hawkes classique} — Trop simpliste mais cadre validé
\end{enumerate}
\end{block}

\vspace{0.5em}
\centering
$\Rightarrow$ Le \textbf{Neural Hawkes Process} est retenu comme base pour la phase suivante
\end{frame}


% ╔═══════════════════════════════════════════════════════════════════════════╗
% ║  PARTIE 3 — SPÉCIALISATION & AMÉLIORATIONS                              ║
% ╚═══════════════════════════════════════════════════════════════════════════╝
\section{Phase 3 — Spécialisation et améliorations}


% ── Roadmap ──────────────────────────────────────────────────────────────────
\begin{frame}{Axes d'amélioration}
\begin{enumerate}
  \item \textbf{Enrichissement des features} : 4 $\to$ 12 features (azimuth, densité, tendance distance, ratio IC\ldots)
  \item \textbf{Architecture Transformer} avec encodage positionnel continu
  \item \textbf{Sortie Gaussienne en log-space} : résout le biais de la régression Huber
  \item \textbf{Estimation d'incertitude} (MC Dropout, variationnel)
  \item \textbf{Modèles spécialisés} par aéroport
  \item \textbf{Spatial Mixture of Experts (MoE)} : $N=3$ experts guidés par la cinématique
\end{enumerate}
\end{frame}


% ── Features enrichies ───────────────────────────────────────────────────────
\begin{frame}{Enrichissement des features}
\small
\begin{tabular}{@{}lll@{}}
\toprule
\textbf{\#} & \textbf{Feature} & \textbf{Description} \\
\midrule
0--3  & icloud, amplitude, dist, maxis & Phase 2 (inchangées) \\
4--5  & azimuth\_sin, azimuth\_cos    & Direction encodée cycliquement \\
6     & dt\_since\_last               & Temps inter-événement \\
7     & dt\_since\_last\_cg           & Temps depuis le dernier CG \\
8     & density\_5min                 & Nb éclairs / 5 min \\
9     & ic\_ratio\_recent             & Ratio IC récent (10 derniers) \\
10    & dist\_to\_centroid            & Distance au centroïde mobile \\
11    & dist\_trend                   & Tendance distance (+ = éloignement) \\
\bottomrule
\end{tabular}

\vspace{0.8em}
\textbf{Motivation :} capturer les signaux de \textbf{dissipation} (éloignement, baisse de densité, ratio IC $\nearrow$)
\end{frame}


% ── Log Gaussian ─────────────────────────────────────────────────────────────
\begin{frame}{Sortie Gaussienne en espace log — la clé}
\begin{columns}[T]
\begin{column}{0.55\textwidth}
\textbf{Problème :} la distribution du temps restant est très asymétrique $\Rightarrow$ la Huber loss biaise vers les grandes valeurs.

\vspace{0.5em}
\textbf{Solution :} $z = \log(1+y)$ normalise la distribution.

$$\mathcal{L}_{\text{Gauss}} = \frac{1}{2}\left[\frac{(z-\mu)^2}{\sigma^2} + 2\log\sigma\right]$$

\vspace{0.5em}
La prédiction finale : $\hat{y} = e^\mu - 1$

L'incertitude est \textbf{naturellement calibrée}.
\end{column}
\begin{column}{0.42\textwidth}
\begin{block}{Impact (la plus grande amélioration)}
\begin{tabular}{@{}lcc@{}}
         & \textbf{Avant} & \textbf{Après} \\
\midrule
Biais    & +11.04 & \hig{+1.96} \\
MAE      & 26.59  & \hig{21.33} \\
Médiane  & 24.18  & \hig{15.96} \\
\end{tabular}

\vspace{0.5em}
\small Biais : $-82\%$\\
MAE : $-20\%$\\
Médiane : $-34\%$
\end{block}
\end{column}
\end{columns}
\end{frame}


% ── Résultats comparatifs ────────────────────────────────────────────────────
\begin{frame}{Résultats comparatifs — toutes variantes}
\centering
\includegraphics[width=0.85\linewidth]{figures/phase2bis_ablation.pdf}

\vspace{0.3em}
\small
\begin{tabular}{llrrrrr}
\toprule
\textbf{It.} & \textbf{Variante} & \textbf{MAE} & \textbf{RMSE} & \textbf{Méd.} & \textbf{Biais} & \textbf{P90} \\
\midrule
Réf. & Baseline 30 min & 60.26 & --- & --- & $-$52.32 & --- \\
Réf. & GRU v1 (Phase 2) & 22.11 & 32.50 & 18.00 & $-$1.61 & 40.12 \\
\rowcolor{meteoGray}\textbf{v3} & \textbf{GRU Gaussian (log)} & \textbf{21.33} & \textbf{30.25} & \textbf{15.96} & \textbf{+1.96} & \textbf{42.13} \\
v3 & Transformer Gaussian & 23.25 & 31.59 & 19.34 & +4.23 & 43.12 \\
\bottomrule
\end{tabular}
\end{frame}




% ── Par aéroport ─────────────────────────────────────────────────────────────
\begin{frame}{Performance par aéroport}
\centering
\includegraphics[width=0.72\linewidth]{figures/phase2bis_per_airport.pdf}

\vspace{0.3em}
\small
Homogénéité entre aéroports : écart max de \textbf{2.57 min} sur la MAE.\\
La spécialisation par aéroport n'améliore que marginalement ($\Delta$ MAE $\approx$ +0.60 min)
\end{frame}




% ╔═══════════════════════════════════════════════════════════════════════════╗
% ║  CONCLUSION                                                              ║
% ╚═══════════════════════════════════════════════════════════════════════════╝
\section{Conclusion et démonstration}


% ── Résultats clés ───────────────────────────────────────────────────────────
\begin{frame}{Résultats finaux}
\begin{columns}[T]
\begin{column}{0.52\textwidth}
\begin{block}{Modèle final — GRU Gaussian (log)}
\begin{tabular}{@{}rl@{}}
MAE          & \hig{21.33 min} (vs 60.26 baseline) \\
RMSE         & 30.25 min \\
Médiane      & 15.96 min \\
Biais        & +1.96 min \\
Calibration  & 96.4\% dans 2$\sigma$ \\
\end{tabular}
\end{block}
\end{column}
\begin{column}{0.45\textwidth}
\begin{block}{Progression}
\centering
\begin{tabular}{@{}lc@{}}
Baseline     & 60.26 min \\
XGBoost AFT  & 57.28 min \\
Neural Hawkes v1 & 22.11 min \\
\textbf{GRU Gaussian} & \textbf{21.33 min} \\
\end{tabular}

\vspace{0.5em}
\small Réduction de \hig{$-65\%$} par rapport à la baseline
\end{block}
\end{column}
\end{columns}

\vspace{0.8em}
\centering
\textbf{+ Interface web interactive} pour explorer les sessions orageuses et les prédictions\\
\small\url{https://mstrdav.github.io/prediction-orage/interactive/}
\end{frame}


% ── Perspectives ─────────────────────────────────────────────────────────────
\begin{frame}{Perspectives}
\begin{enumerate}
  \item \textbf{Plus de données} : le dataset (1\,682 sessions) reste limité
  \item \textbf{Features météo} : données radar (réflectivité, vent)
  \item \textbf{Recalibration} isotonique pour corriger la sur-couverture
  \item \textbf{Hybride MoE + GRU} : combiner la puissance séquentielle du GRU avec la spécialisation spatiale du MoE
  \item \textbf{Déploiement temps réel} : intégration dans les systèmes d'alerte existants
\end{enumerate}
\end{frame}


% ── Merci ─────────────────────────────────────────────────────────────────────
\begin{frame}[standout]
Merci pour votre attention !

\vspace{1em}
\large Questions ?

\vspace{1.5em}
\normalsize
\textbf{Les Vieux Briscards}\\
\small Colin Geindre · Alexandre Hayoun · Tao Levray

\vspace{1.5em}
\footnotesize
\url{https://github.com/Mstrdav/prediction-orage}\\
\url{https://mstrdav.github.io/prediction-orage/interactive/}
\end{frame}

\end{document}
